\documentclass[11pt]{article}
\usepackage[T1]{fontenc}
\usepackage{lmodern,amsmath,amssymb,amsthm,microtype}
\usepackage[margin=1in]{geometry}
\usepackage[hidelinks]{hyperref}
\newtheorem{theorem}{Theorem}
\newtheorem{lemma}[theorem]{Lemma}
\newcommand{\Z}{\mathbb Z}
\newcommand{\R}{\mathbb R}
\newcommand{\F}{\mathbb F}
\newcommand{\code}[1]{\texttt{\detokenize{#1}}}
\hypersetup{pdftitle={Every n-point real set contains a Sidon subset of size at least (2/(3 sqrt(3)) - o(1)) sqrt(n)},pdfauthor={Dohyun Park}}
\title{Every $n$-point real set contains a Sidon subset\protect\\ of size at least $\bigl(2/(3\sqrt{3})-o(1)\bigr)\sqrt{n}$}
\author{Dohyun Park\\[3pt]
\small Department of Mechanical Engineering, Seoul National University\\
\small\href{mailto:eddiepark92@gmail.com}{\texttt{eddiepark92@gmail.com}}}
\date{November 25, 2026}
\begin{document}
\maketitle

\begin{abstract}
For every $\varepsilon>0$ and all sufficiently large $n$, every set of $n$ real numbers contains a Sidon subset of size at least $(2/(3\sqrt{3})-\varepsilon)\sqrt{n}$. The coefficient doubles the one obtained by Bailleul and Riblet and is $\sqrt{2}$ times the coefficient reported in an August 2026 SciNet finding. The conjectured coefficient $1$ remains open. The proof samples a random affine map into a prime field and counts occupied residue classes indexed by a cyclic Sidon set. A finite version guarantees at least $385$ elements when $n=10^6$. The main theorem has been formalized in Lean.
\end{abstract}

\section{The guaranteed size}

A subset \(S\) of an abelian group is \emph{Sidon} if
\(a+b=c+d\), with \(a,b,c,d\in S\), implies
\((a,b)=(c,d)\) or \((a,b)=(d,c)\). The definition includes repeated
summands. Write \(H(n)\) for the minimum, over \(n\)-element subsets of
\(\R\), of the size of a largest Sidon subset. The problem database credits the original question to Riddell \cite{riddell}. In the formulation of Erd\H{o}s problem~\#530, the question is whether
\(H(n)=(1+o(1))\sqrt n\) \cite{erdos}.

\begin{theorem}\label{th:optimal}
For every $\varepsilon>0$ there is an $n_0$ such that every finite $A\subset\R$ of size $n\ge n_0$ contains a Sidon subset of cardinality at least
\[
 \left(\frac{2}{3\sqrt3}-\varepsilon\right)\sqrt n .
\]
Equivalently, $\displaystyle\liminf_{n\to\infty} H(n)/\sqrt n\ge 2/(3\sqrt3)$.
\end{theorem}

Koml\'os, Sulyok and Szemer\'edi established a general comparison theorem
for sets avoiding linear configurations \cite{KSS}. For Sidon subsets,
Abbott obtained every coefficient below \(2/25\) \cite{abbott}.
Bailleul and Riblet proved a lower coefficient of
\(1/(3\sqrt3)=0.19245\ldots\) by compressing a subset into a cyclic
group and applying a Singer covering \cite{BR}. An August 2026
agent-authored SciNet posting reports an improvement to
\(\sqrt6/9=0.27216\ldots\). It uses reflection to halve the pair-collision
estimate in the unshifted dilation argument \cite{scinet}.
The coefficient in Theorem~\ref{th:optimal} is twice the Bailleul--Riblet coefficient for real sets and $\sqrt2$ times the coefficient in that public posting. I am not aware of a better universal lower coefficient for arbitrary finite sets of real numbers. The
coefficient \(1\) in the Erd\H{o}s question remains open.

\begin{theorem}\label{th:main}
For every integer \(n\ge0\),
\[
  36n\le343\bigl(H(n)+4\bigr)^2,\qquad
  H(n)\ge\left\lceil \frac{6}{7\sqrt7}\sqrt n-4\right\rceil .
\]
\end{theorem}

\begin{theorem}\label{th:param}
For each prime \(q\), put \(m=q(q-1)\) and \(b=q-1\). For every \(n\ge0\),
\begin{equation}\label{eq:param}
  2mbn\ \le\ 4m^2 H(n)+b\binom n2 .
\end{equation}
If \(q\ge3\) and \(n_q=\lfloor4q(q-1)/3\rfloor\), then
\[
  H(n_q)\ \ge\ \left\lceil\frac{4(q-1)}9\right\rceil .
\]
In particular,
\(\displaystyle\limsup_{n\to\infty}H(n)/\sqrt n\ge2/(3\sqrt3)\).
\end{theorem}

For \(n=10^6\), Theorem~\ref{th:main} guarantees \(320\) elements.
The parameter inequality gives more information at this size. Taking
the prime \(q=863\) in \eqref{eq:param} gives \(H(10^6)\ge385\),
by rounding a lower bound of \(384.6676\ldots\) upwards.

\section{Occupancy of a cyclic Sidon set}

We first work over the integers. Let \(A\subset\Z\) have \(n\) elements
and let \(B\subset\Z/m\Z\) be Sidon with \(|B|=b\) and \(m\ge1\).
Write \(h(A)\) for the maximum size of a Sidon subset of \(A\).

\begin{lemma}\label{lem:template}
For every such \(A,B\),
\[
  2mbn\le 4m^2 h(A)+b\binom n2 .
\]
\end{lemma}
\begin{proof}
Take a prime \(P\) larger than every absolute difference between elements
of \(A\). Choose \((u,v)\) uniformly from \(\F_P^2\), allowing \(u=0\),
and write \(X_a\) for the representative in \([0,P)\) of \(ua+v\).
For \(a\ne a'\), the map \((u,v)\mapsto(ua+v,ua'+v)\) is bijective.
The difference of the two prescribed values determines \(u\), then either
value determines \(v\). Thus each \(X_a\) is uniform and each pair
\((X_a,X_{a'})\) has the uniform joint distribution.

Set \(L=\lfloor P/(2m)\rfloor\). Represent \(t\in B\) by its integer
in \([0,m)\) and associate to it the window
\[
 W_t=\{t+mj:0\le j<L\}.
\]
Each window has \(L\) residues, all strictly less than \(mL\le P/2\).
The probability that a fixed \(a\) hits \(W_t\) is \(L/P\). Two distinct
points both hit \(W_t\) with probability \(L^2/P^2\). If \(K_t\) is the
number of inputs in \(W_t\), then
\(\mathbf1_{\{K_t>0\}}\ge K_t-\binom{K_t}{2}\) for each integer
\(K_t\ge0\). Taking expectations and summing over \(t\in B\), some
choice of \((u,v)\) occupies at least
\begin{equation}\label{eq:occ}
  b\left(\frac{nL}{P}-\binom n2\frac{L^2}{P^2}\right)
\end{equation}
distinct windows. This argument uses only the distribution of one or two
images at a time.

Choose one input from each occupied window to form \(S\).
Suppose \(a+b'=c+d\) among elements of \(S\). The two sums of their four
affine images are congruent modulo \(P\). Both sums of their nonnegative
representatives are strictly less than \(P\), so they are equal as
integers. Reduction modulo \(m\) gives an equality of pair sums among the
window labels. Since \(B\) is cyclic Sidon, the unordered pairs of labels
coincide. Each label has exactly one chosen representative in \(S\), so
the unordered pairs of inputs coincide as well.

Consequently \(h(A)\) bounds \eqref{eq:occ} for every sufficiently
large prime \(P\). Letting \(P\) tend to infinity gives the stated inequality.
The Lean proof removes the floor by integer arithmetic. For \(D=2m\)
it takes \(P>D^2bn\) as well as injectivity on \(A\). The inequalities
\(DL\le P\le DL+D\) then give the same integral conclusion.
\end{proof}

For the modulus \(m=q(q-1)\), Ruzsa's construction
\cite[Theorem~4.4]{ruzsa} gives \(b=q-1\). To see this, choose a
generator \(g\) of \(\F_q^\times\). Under the Chinese remainder
identification, take the graph
\[
  B=\{(i,g^i):i\in\Z/(q-1)\Z\}
     \subset \Z/(q-1)\Z\times\F_q .
\]
Equality of two graph pair sums gives equality of the sums of the
second coordinates. Equality of the first-coordinate sums gives equality
of their products. The unordered pair of field elements is recovered as
the multiset of roots of its monic quadratic. The distinct powers of
\(g\) recover the indices. Lemma~\ref{lem:template} therefore proves
\eqref{eq:param} for integer inputs.

For a finite real set \(A\), take a rational linear functional on the
rational span of \(A\) that is injective on \(A\cup(A+A)\). Such a
functional exists because the finitely many proper hyperplanes in the
dual corresponding to differences of distinct points in this union cannot
cover the dual over \(\mathbb Q\). Clearing the denominators of the
images of \(A\) produces an injective integer model preserving and
reflecting every equality of two-term sums. Any Sidon subset in the model
thus pulls back with its full cardinality, proving \eqref{eq:param}
for real inputs.

\section{Choosing the prime}

\begin{proof}[Proof of Theorem~\ref{th:main}]
Let \(r=\lfloor\sqrt{n/112}\rfloor\).
For \(r=0\) the assertion follows even from the empty subset.
For \(r\ge1\), put \(k=7r\) and restrict attention to
\(t=16r(7r-1)\le112r^2\le n\) of the input elements.
Bertrand's postulate gives a prime \(q\) with \(k<q\le2k\).
Writing \(y=k/q\in[1/2,1]\), \eqref{eq:param} gives a subset of size
\(s\) satisfying
\begin{align*}
s &\ge \frac{t}{2q}\left(1-\frac{t-1}{4q(q-1)}\right)\\
  &\ge \frac87(k-1)y\left(1-\frac47 y^2\right)
   \ge \frac{24}{49}(k-1).
\end{align*}
Indeed, \(t=16k(k-1)/7\) and \((k-1)/(q-1)\le k/q\).
The last inequality follows from
\(7y-4y^3-3=(2y-1)(1-y)(2y+3)\ge0\).
Thus \(7(s+4)\ge24(r+1)\). Squaring and using
\(n<112(r+1)^2\) proves \(36n\le343(s+4)^2\).
\end{proof}

\begin{proof}[Proof of the prime-indexed conclusion in Theorem~\ref{th:param}]
Let \(m=q(q-1)\), \(n=\lfloor4m/3\rfloor\) and \(e=4m-3n\in\{0,1,2\}\).
Then
\[
  36mn-9n(n-1)-32m^2=12m-4me-3e-e^2\ge0,
\]
since \(m\ge6\). Substituting in \eqref{eq:param} gives
\(9H(n)\ge4(q-1)\). As the primes \(q\) increase, \(n\sim4q^2/3\)
and \(4(q-1)/(9\sqrt n)\to2/(3\sqrt3)\).
\end{proof}

\begin{proof}[Proof of Theorem~\ref{th:optimal}]
For each $n\ge1$ choose the least prime $q=q(n)$ exceeding $(\sqrt3/2)\sqrt n$. For each fixed $\eta>0$ there is a prime between $x$ and $(1+\eta)x$ for all sufficiently large $x$, by the prime number theorem. Thus $q/\sqrt n\to z=\sqrt3/2$ and $(q-1)/\sqrt n\to z$.
Dividing \eqref{eq:param} by $4m^2\sqrt n$ gives
\[
 \frac{H(n)}{\sqrt n}\ \ge\ \frac{1}{2(q/\sqrt n)}
 -\frac{1-1/n}{8((q-1)/\sqrt n)(q/\sqrt n)^2}.
\]
The right side tends to $1/(2z)-1/(8z^3)=2/(3\sqrt3)$. The choice of $q(n)$ depends only on the input size, so this estimate holds uniformly over all $A$. The value $z=\sqrt3/2$ maximizes the limiting expression for this estimate.
\end{proof}

\section{Attribution and verification}

Pairwise-uniform affine evaluations and restriction to a half cube appear
in Green and Wolf's torus construction for progression-free integer sets
\cite[Lemma~2.1]{GW}. O'Bryant used a common dilation and translation
for subsets of arbitrary integer sets avoiding progressions \cite{obryant}.
Kohayakawa, Lee, R\"odl and Samotij selected representatives of occupied
intervals indexed by a Sidon set in a random ambient set
\cite[Lemma~7.8]{KLRS}. Pach and Zakharov used a finite integer model
and random dilation into a half cube in their work on bi-Sidon sets
\cite[Propositions~1--2]{PZ}.
Here the independent additive shift makes the two-hit probability exactly
\((L/P)^2\) for each label, and the deterministic pullback makes the
result apply to every finite input.

\paragraph{Formalization.}
The finite estimates are formalized in the namespace \code{SidonCompression}. The theorem \code{real_sidon_bound} proves
\eqref{eq:param}. The declaration \code{real_sidon_peak} gives the finite bound at $n_q$,
and \code{strong_bertrand_sidon_real} proves Theorem~\ref{th:main}.
The multiplicative prime-interval estimate used in Theorem~\ref{th:optimal}
is \code{prime_between} from the pinned \emph{Prime Number Theorem and More}
development \cite{pnt}.

\section*{Acknowledgements and disclosure}
The proof and the Lean formalization were produced with assistance from
ChatGPT (OpenAI) and Claude (Anthropic). The main theorem,
\code{SidonCompression.real_sidon_asymptotic}, was machine-checked in
Lean 4.19.0 with Mathlib at revision \code{c44e0c8ee63c}.
An independent rebuild reproduced the kernel's transitive axiom report.
For both the main theorem and \code{prime_between}, the complete list is
\code{propext}, \code{Classical.choice}, and \code{Quot.sound}.

\begin{thebibliography}{99}\small\setlength{\itemsep}{1pt}
\bibitem{abbott} H.\,L. Abbott. Sidon sets. \emph{Canadian Mathematical
Bulletin} \textbf{33} (1990), 335--341.
\href{https://doi.org/10.4153/CMB-1990-056-6}{doi:10.4153/CMB-1990-056-6}.
\bibitem{BR} A. Bailleul and R. Riblet. On the largest Sidon subset in a
finite subset of \(\mathbb R^N\). \emph{European Journal of Combinatorics} \textbf{139} (2027),
article 104443. \href{https://doi.org/10.1016/j.ejc.2026.104443}{Journal record} and
\href{https://arxiv.org/abs/2605.03181}{arXiv:2605.03181}.
\bibitem{erdos} T. Bloom, ed. Erd\H{o}s Problems, Problem \#530.
\href{https://www.erdosproblems.com/530}{Problem statement} and
\href{https://www.erdosproblems.com/forum/thread/530}{discussion}.
\bibitem{GW} B. Green and J. Wolf. A note on Elkin's improvement of
Behrend's construction. In D. Chudnovsky and G. Chudnovsky (eds.), \emph{Additive Number Theory}, Springer, 2010, 141--144. \href{https://arxiv.org/abs/0810.0732}{arXiv:0810.0732}.
\bibitem{KLRS} Y. Kohayakawa, S.\,J. Lee, V. R\"odl and W. Samotij.
The number of Sidon sets and the maximum size of Sidon sets contained in
a sparse random set of integers. \emph{Random Structures \& Algorithms}
\textbf{46} (2015), 1--25. \href{https://doi.org/10.1002/rsa.20496}{doi:10.1002/rsa.20496}.
\bibitem{KSS} J. Koml\'os, M. Sulyok and E. Szemer\'edi.
Linear problems in combinatorial number theory. \emph{Acta Mathematica
Academiae Scientiarum Hungaricae} \textbf{26} (1975), 113--121.
\href{https://doi.org/10.1007/BF01895954}{doi:10.1007/BF01895954}.
\bibitem{obryant} K. O'Bryant. Thick subsets that do not contain
arithmetic progressions. \emph{Integers} \textbf{13} (2013), A18.
\href{https://math.colgate.edu/~integers/n18/n18.pdf}{Full text}.
\bibitem{PZ} J. Pach and D. Zakharov. Ruzsa's problem on bi-Sidon sets.
\emph{Combinatorica} \textbf{45} (2025), article 26.
\href{https://doi.org/10.1007/s00493-025-00151-5}{doi:10.1007/s00493-025-00151-5}.
\bibitem{riddell} J. Riddell. On sets of numbers containing no $l$ terms in arithmetic progression.
\emph{Nieuw Archief voor Wiskunde} (3) \textbf{17} (1969), 204--209.
\bibitem{ruzsa} I.\,Z. Ruzsa. Solving a linear equation in a set of
integers I. \emph{Acta Arithmetica} \textbf{65} (1993), 259--282.
\href{https://doi.org/10.4064/aa-65-3-259-282}{doi:10.4064/aa-65-3-259-282}.
\bibitem{scinet} Proof-Track Strategist. Improved constant for guaranteed
Sidon subsets (Erd\H{o}s \#530): $\sqrt{3}/9\to\sqrt{6}/9$, sharpening one estimate in Bailleul--Riblet. Public agent-authored finding, SciNet,
2 August 2026. \href{https://api.scinet.pub/f/6e6583f0-89cb-4c69-9f76-a7e5ed0da6e8}{Finding 6e6583f0}.
\bibitem{pnt}
A. Kontorovich, T. Tao and contributors, \emph{Prime Number Theorem and More}, Lean formalization, \href{https://github.com/AlexKontorovich/PrimeNumberTheoremAnd/tree/426b540a86d5b4a103cec5528e68168ea172da66}{revision 426b540a86d5}, \code{PrimeNumberTheoremAnd.Consequences}.
\end{thebibliography}
\end{document}
